\clearpage
\section*{Developer funding in the main specifications}
\textit{September 14, 2026. Agent-drafted exploratory results on score\_robustness; submitted scores and manuscript unchanged.}

Ordering aldermen by developer funding produces imprecise density and price differences, but the permitting result remains negative and statistically significant. The signed permit estimate is $-10.7\%$ ($p=0.038$, 95\% interval $[-19.8,-0.7]\%$). Higher funding means a larger share of eligible campaign dollars identified as developer-linked. It is not labelled greater stringency or leniency, and the direction is held fixed throughout this comparison.

For all construction, the boundary estimates are $-2.5\%$ for FAR ($p=0.620$) and $-1.2\%$ for DUPAC ($p=0.868$). For multifamily construction they are $+2.9\%$ for FAR ($p=0.748$) and $+17.6\%$ for DUPAC ($p=0.245$). Rents differ by $-2.3\%$ ($p=0.168$), and sale prices by $+3.0\%$ ($p=0.143$). These estimates compare the nearest 100 feet on either side using the paper's ten-bin specification within 500 feet. The wide intervals do not establish zero effects. In particular, the sale-price magnitude remains near the original estimate while becoming less precise.

The audit reestimates each original specification on its full sample and then on the observations usable under the developer ranking. It reproduces all seven original coefficients and standard errors to the paper's displayed three-decimal precision. Both common-sample versions retain exactly the same fitted observations, checked directly. The full-sample permit estimate is $-11.9\%$; the original ranking on the developer-supported observations gives $-8.0\%$, compared with $-10.7\%$ under developer ordering. Thus the comparison changes both sample coverage and ordering, and these effects are reported separately. Developer ordering reverses 166 of the 583 reassigned blocks with supported comparisons. The developer event study's joint pretrend test has $p=0.500$; this does not establish all identifying assumptions. Separate directional estimates are $-14.8\%$ toward higher developer funding ($p=0.113$) and $+7.6\%$ toward lower funding ($p=0.238$). The significant signed result imposes equal-and-opposite log effects, as in the main permit specification.

The score uses the prior audit's documented developer definition and personal-campaign receipt restrictions, with no new donor classifications or weights chosen to improve agreement. Boundary scores use in-office receipts during 2006--2022. The permit-remap score uses only 2006--2014 receipts, before the 2015 reassignment. An observed zero developer share is retained; missing fundraising is not imputed as zero. Missing or tied scores prevent ordering. In the permit analysis the entire unsupported ward pair is excluded, including its unchanged controls, so a tied reassignment is not treated as an unchanged block. Common fitted samples contain 3,480 construction projects (794 multifamily), 230,189 rent observations, 42,097 sales, and 31,487 permit block-years.

The pattern does not establish that campaign donations measure regulatory behavior. Many donor firms remain unclassified, the boundary score is retrospective, and the two measures may capture different aspects of local politics. The developer-permit rank correlation was only 0.163 in the prior audit. This exercise is exploratory and does not change the submitted estimates.

\textit{Reproduction.} Run \texttt{make} in
\path{tasks/audits/developer_score_results/code/} and open
\path{../output/developer_results.html}. Two scripts estimate and display the comparisons. Four production analysis files were copied unchanged to \path{data_raw/score_robustness/}; \path{sources/benchmark_sources.csv} records their producers and hashes. Donation receipts come from the preceding audit at \texttt{bc100d9b}. These benchmark files are preserved locally.
